\documentclass[10pt,a4paper]{amsart}
\usepackage[margin=19mm]{geometry}
\usepackage{amsmath,amssymb,mathtools}
\usepackage[scr=rsfs,cal=euler]{mathalfa}
\usepackage{xfrac}
\usepackage[unicode,hidelinks,bookmarks=false]{hyperref}
\newcommand{\ie}{i.e.\ }
\newcommand{\cf}{cf.\ }
\newcommand{\resp}{respectively\ }
\newcommand{\Subsection}{\subsection}
\providecommand{\nobreakdash}{\nobreak\hbox{-}}
\setlength{\emergencystretch}{2em}
\multlinegap0pt
\usepackage{bm}
\usepackage{bbm}
\usepackage{dsfont}
\usepackage{xspace}
\usepackage{cancel}
\usepackage{mathtools}
\usepackage[shortlabels]{enumitem}
\usepackage{amsthm}
\makeatletter
\@ifundefined{thm}{\newtheorem{thm}{Thm}}{}
\@ifundefined{lem}{\newtheorem{lem}[thm]{Lemma}}{}
\makeatother
\newcommand{\ud}{\, \mathrm{d}}
\title[Confinement transitions in half-space constrained Riesz gase]{Confinement transitions in half-space constrained Riesz gases}
\author{Sung-Soo Byun, Yong-Woo Lee, Eui Yoo}
\date{Source: arXiv:2608.11813v1 (CC BY 4.0)}
\begin{document}
\maketitle

\noindent\emph{Source and licence.} Confinement transitions in half-space constrained Riesz gases. Version arXiv:2608.11813v1 (CC BY 4.0), distributed under the Creative Commons Attribution 4.0 International licence, as recorded in the arXiv licence metadata for this exact version and shown on the abstract page https://arxiv.org/abs/2608.11813v1. Full rights notice: licenses/arxiv-2608.11813-CC-BY-4.0.txt.\par\smallskip

\noindent\textit{Editorial scope.} Counted unit: the Lem whose source label is \texttt{lem: property f(x)} in the article (source0.tex lines 1174--1203, source label \texttt{lem: property f(x)}), delivered together with the article's own complete original proof of it (source0.tex lines 1216--1333, one \texttt{proof} environment of 118 lines). No mathematics is written, repaired or re-derived: statement and proof are the article's own text line for line. The proof is detached from the statement in the source: the article states the result at lines 1174--1203 and prints its proof at lines 1216--1333, and the proof environment's own header names the statement, so the association is the article's own. Required context retained uncounted: eq:critical point incomplete beta (lines 804--821); def:f(x) (lines 1153--1155); eq:f'(x) (lines 1159--1165). Every definition, display and label of those lines is retained unchanged. Omitted from this package and not redistributed: 1--803, 822--1152, 1156--1158, 1166--1173, 1204--1215, 1334--1903. Nothing inside a retained excerpt is omitted or altered: every retained body is delivered exactly as the article's lines are written, so no omission receipt of any kind accompanies this package. Citations: the retained excerpts carry no citation command at all, so no bibliography entry of the article is transcribed and no reference is redistributed. Reference adaptation: none was needed and none was made; every reference inside the delivered unit resolves inside it, so no unresolved reference or citation is produced. Printed slips kept literal: no printed word, symbol, hypothesis or number of the counted statement or of its proof was altered. Layout and class: the article's own class is \texttt{amsart} with packages including amssymb, amsthm, amsfonts, amsmath, bbm, pmboxdraw, verbatim, graphicx, color, hyperref, cite, ulem, subcaption, tabulary; the wrapper is the lane's pinned \texttt{amsart} editorial wrapper at 10pt on A4, which restates the theorem-like environments the retained text needs and adds the article's own macro definitions that the retained text needs (\texttt{\textbackslash ud}), with extra wrapper packages (bm, bbm, dsfont, xspace, cancel, mathtools, enumitem). The delivered numbering is the unit's own. Assets: no retained span contains a figure environment, a graphics-inclusion command, a PDF-page inclusion, a table or a TikZ picture; no image file is copied into this package, the package declares no asset dependency and no image file is redistributed with it. Licence: version arXiv:2608.11813v1 of this article is distributed under the Creative Commons Attribution 4.0 International licence, as recorded in the arXiv licence metadata for this exact version and shown on the abstract page https://arxiv.org/abs/2608.11813v1. A full rights notice is delivered at licenses/arxiv-2608.11813-CC-BY-4.0.txt. Characters: every retained excerpt is pure ASCII, so the delivered engine, XeLaTeX with its default Latin Modern text font, reports no missing character.
\begin{equation}
\label{eq:critical point incomplete beta}
\mathsf{x}^{-2\alpha}
    (1+\mathsf{x}^2)^{\frac{3-d}{2}}
    +
    \Big(
        \frac{s}{2}
        -
        \frac{\alpha}{\mathsf{x}^2}
    \Big)
    \mathrm{B}_{\frac{\mathsf{x}^2}{1+\mathsf{x}^2}}
    \Big(
        \frac{d-s-1}{2},
        \frac{s}{2}
    \Big)
    =
    \mathsf{C}_{d-1,s},
\end{equation}

\begin{equation}\label{def:f(x)}
    f(x) := x^{1-2\alpha}(1+x^2)^{-\frac{d-3}{2}} +\Big(\frac{s}{2}x +\frac{\alpha}{x}\Big) \mathrm{B}_{\frac{x^2}{1+x^2}}\Big(\frac{d-s-1}{2}, \frac{s}{2}\Big)-\mathsf{C}_{d-1, s}x, 
\end{equation}

\begin{align}
\label{eq:f'(x)}
\begin{split}
    f'(x) 
    &= x^{-2\alpha}(1+x^2)^{-\frac{d-3}{2}} + \Big(\frac{s}{2}-\frac{\alpha}{x^2}\Big)\mathrm{B}_{\frac{x^2}{1+x^2}}\Big(\frac{d-s-1}{2},\frac{s}{2}\Big)-\mathsf{C}_{d-1,s}.
\end{split}
\end{align}

\begin{lem} \label{lem: property f(x)}
Assume that $s\in(d-3,d-1)$. The function $f$ defined in~\eqref{def:f(x)} satisfies the following properties: 
\begin{enumerate}[label=\textup{(\roman*)}]
    \item \label{item: limits f(x)}
  As $x\to\infty$, we have $f(x)\to-\infty$, whereas 
    \begin{equation}
         \lim_{x\to 0^+}f(x)=  
        \begin{cases}
            0, &\textup{if } s \in (d-3,d-2),
            \smallskip 
            \\
            2, &\textup{if } s= d-2,
            \smallskip 
            \\
            +\infty, &\textup{if } s\in (d-2,d-1).
        \end{cases}
    \end{equation}
    Consequently, $f$ attains a global maximum on $[0,\infty)$ for $s\in(d-3,d-2]$, whereas no finite global maximum exists for $s\in(d-2,d-1)$. 
    \smallskip 
    \item \label{item: maximum existence}
    In the strongly long-ranged regime $s\in(d-3,d-2]$, the function $f$ attains its unique global maximum at $x=\mathsf{x}$, where $\mathsf{x}$ is given by~\eqref{eq:critical point incomplete beta} for $s\in (d-3,d-2)$ and $\mathsf{x}=0$ for $s=d-2$.
    Moreover, the maximum value is given by
    \begin{equation}
    \label{eq:f(x) maximum characterisation}
        \max_{x\in [0,\infty)} f(x) =f(\mathsf{x})= \frac{2\alpha}{\mathsf{x}} \mathrm{B}_{\frac{\mathsf{x}^2}{1+\mathsf{x}^2}}\Big(\frac{d-s-1}{2},\frac{s}{2}\Big),
    \end{equation}
   where, in the case $\mathsf{x}=0$, the right-hand side is understood in
the limiting sense as $\mathsf{x}\to0^+$. 
\end{enumerate}
\end{lem}

\begin{proof}[Proof of Lemma~\ref{lem: property f(x)}]
We first show \textrm{(i)}. Note that as $x\to\infty$, we have 
\begin{equation}\label{eq:Beta completion}
    \lim_{x\to \infty}\Big(\frac{s}{2}\mathrm{B}_{\frac{x^2}{1+x^2}}\Big(\frac{d-s-1}{2},\frac{s}{2}\Big)+x^{-s}\Big) 
    = \frac{\Gamma(\frac{d-s-1}{2})\Gamma(\frac{s}{2}+1)} {\Gamma(\frac{d-1}{2})}
    = \mathsf{C}_{d-1,s}\frac{d-1}{s+2},
\end{equation}
It therefore follows from~\eqref{def:f(x)} that, for all $s\in(d-3,d-1)$,
\begin{equation}
    \lim_{x\to\infty} \frac{f(x)}{x} =\lim_{x\to \infty}
    \Big(\frac{s}{2}\mathrm{B}_{\frac{x^2}{1+x^2}}\Big(\frac{d-s-1}{2},\frac{s}{2}\Big)+x^{-s}\Big) - \mathsf{C}_{d-1,s} = \mathsf{C}_{d-1,s} \frac{(d-3)-s}{s+2} <0.
\end{equation}
Consequently, $f(x)\to-\infty$ as $x\to\infty$. 

For the limit as $x\to0^+$, we first observe that 
\begin{equation}\label{eq:Beta to zero}
    \mathrm{B}_{\frac{x^2}{1+x^2}}\Big(\frac{d-s-1}{2},\frac{s}{2}\Big) =\int_0^{\frac{x^2}{1+x^2}} u^{-\alpha} (1-u)^{\frac{s}{2}-1}\ud u = \frac{x^{2-2\alpha}}{1-\alpha}\big(1+ O(x^2)\big), \quad x\to 0^+.
\end{equation}
For $s\in(d-3,d-2)$, we have $\alpha\in(0,\frac12)$, and hence \eqref{eq:Beta to zero} together with~\eqref{def:f(x)} readily gives $\lim_{x\to 0^+}f(x)=0$. On the other hand, for $s\in[d-2,d-1)$, we have
\begin{align*}
    \lim_{x\to0^+}\frac{f(x)}{x^{1-2\alpha}}
    &=
    \lim_{x\to0^+}
    \bigg[
        (1+x^2)^{-\frac{d-3}{2}}
        +
        \Big(\frac{s}{2}x+\frac{\alpha}{x}\Big)
        \frac{1}{x^{1-2\alpha}}
        \mathrm{B}_{\frac{x^2}{1+x^2}}
        \Big(
            \frac{d-s-1}{2},\frac{s}{2}
        \Big)
        -
        \mathsf{C}_{d-1,s}x^{2\alpha}
    \bigg]
    =
    \frac{1}{1-\alpha}.
\end{align*} 
Since $\alpha=\frac{1}{2}$ when $s=d-2$, while $\alpha\in(\frac{1}{2},1)$ when $s\in(d-2,d-1)$, the desired limits follow, completing the proof of~\ref{item: limits f(x)}.  

\medskip
Next, we prove~\textrm{(ii)}. We first note the useful representation 
\begin{equation}\label{eq:incomplete Beta integral}
    \mathrm{B}_{\frac{x^2}{1+x^2}}\Big(\frac{d-s-1}{2},\frac{s}{2}\Big) = \int_0^{x^2} v^{-\alpha}(1+v)^{-\frac{d-1}{2}}\ud v= x^{2-2\alpha}\int_0^1 u^{-\alpha}(1+x^2u)^{-\frac{d-1}{2}}\ud u.
\end{equation}
Using~\eqref{eq:incomplete Beta integral}, for $x>0$, we have 
\begin{align}\label{eq:f''(x)}
\begin{split}
    f''(x) 
    &= -2\alpha x^{-1-2\alpha}(1+x^2)^{-\frac{d-1}{2}}\bigg(2- \int_0^1  u^{-\alpha}\Big(\frac{1+x^2}{1+x^2u}\Big)^{\frac{d-1}{2}}\ud u\bigg).
\end{split}
\end{align}
Notice that the function
\[
h(x):= \int_0^1  u^{-\alpha}\Big(\frac{1+x^2}{1+x^2u}\Big)^{\frac{d-1}{2}}\ud u, \quad x\ge 0,
\]
is an increasing function in $x$. In particular, we have
\begin{equation}
    h(0) = \frac{1}{1-\alpha}\le 2,
\end{equation}
where the equality holds only if $s=d-2$. 

We first prove~\ref{item: maximum existence} in the case $s=d-2$. Since
$h(x)\ge h(0)=2$, it follows from~\eqref{eq:f''(x)} that 
\[
    f''(x) = -2\alpha x^{-1-2\alpha}(1+x^2)^{-\frac{d-1}{2}} (2-h(x) )\ge  0, \qquad x>0. 
\]
Thus, $f$ is convex on $(0,\infty)$, and hence $f'$ is increasing. On the
other hand, by~\eqref{eq:f'(x)} and~\eqref{eq:Beta completion}, 
\begin{align*}
    \lim_{x\to\infty} f'(x) &= \lim_{x\to\infty}\frac{s}{2}\mathrm{B}_{\frac{x^2}{1+x^2}}\Big(\frac{1}{2},\frac{s}{2}\Big) -\mathsf{C}_{d-1,d-2}=-\frac{1}{d}\mathsf{C}_{d-1,d-2}<0.
\end{align*}
Since $f'$ is increasing, we conclude that $f'(x)<0$ for all $x>0$.
Therefore, $f$ is strictly decreasing on $(0,\infty)$ and attains its
unique global maximum at $x=\mathsf{x}=0$. Using~\eqref{eq:Beta to zero}, we also obtain 
\[
    \lim_{x\to0^+}\frac{1}{x}\mathrm{B}_{\frac{x^2}{1+x^2}}\Big(\frac{1}{2},\frac{s}{2}\Big) = 2 = f(0),
\]
which shows (ii) for $s=d-2$. 
 

We now turn to the case $s\in(d-3,d-2)$. By~\eqref{eq:f''(x)},
the inequality $h(0)<2$ implies that $f''(x)<0$ for all sufficiently
small $x>0$. Moreover, since $h$ is increasing, $f''$ has at most one
zero on $(0,\infty)$.

We claim that $f$ has a unique critical point on $(0,\infty)$. In view
of~\ref{item: limits f(x)} and the observation following~\eqref{eq:f'(x)},
this claim implies that~\eqref{eq:critical point incomplete beta} admits
a unique solution $\mathsf{x}\in(0,\infty)$ and that $f$ attains its
unique global maximum at $x=\mathsf{x}$.

We argue by contradiction. Since $f''$ has at most one zero on
$(0,\infty)$, the function $f$ has at most two critical points there.
Suppose that $f$ has exactly two critical points, say
$\mathsf{x}_1$ and $\mathsf{x}_2$, with $0<\mathsf{x}_1<\mathsf{x}_2$. Then it follows that $f''$ must vanish at
some point in $(\mathsf{x}_1,\mathsf{x}_2)$. Since $f''(x)<0$ for all
sufficiently small $x>0$ and $f''$ has at most one zero, it follows that
\[
    f''(\mathsf{x}_1)<0, \qquad f''(\mathsf{x}_2)>0.
\]
On the other hand, since $f''(x)>0$ for $x>\mathsf{x}_2$, the function
$f'$ is increasing on $(\mathsf{x}_2,\infty)$. Hence, by
\eqref{eq:f'(x)} and~\eqref{eq:Beta completion}, 
\[
    0=f'(\mathsf{x}_2)\le \lim_{x\to\infty}f'(x) = \lim_{x\to\infty}\frac{s}{2}\mathrm{B}_{\frac{x^2}{1+x^2}}\Big(\frac{d-s-1}{2},\frac{s}{2}\Big)-\mathsf{C}_{d-1,s} = \mathsf{C}_{d-1,s}\frac{(d-3)-s}{s+2}<0,
\]
which is a contradiction. Therefore, $f$ has a unique critical point
on $(0,\infty)$.  

It remains to verify~\eqref{eq:f(x) maximum characterisation} for
$s\in(d-3,d-2)$. From~\eqref{def:f(x)} and~\eqref{eq:f'(x)}, we readily obtain 
\begin{equation}
    f(x)-xf'(x)= \frac{2\alpha}{x}\mathrm{B}_{\frac{x^2}{1+x^2}}\Big(\frac{d-s-1}{2},\frac{s}{2}\Big).
\end{equation}
Evaluating this identity at the unique critical point $x=\mathsf{x}$,
for which $f'(\mathsf{x})=0$, gives the desired result \eqref{eq:f(x) maximum characterisation}. 
\end{proof}

\end{document}
